%% kapa: exam=1 round=24 year=2013 period=1 form=A subject=economics subject_ko=경제원론 seq=13 number=53
%% source: Q-Net 감정평가사 기출문제 (cst00309) artlSeq=5073802 file=extracted/1008/5073802/2013년 제24회 감정평가사 1차시험 기출문제/1교시(A형).hwp
%% conversion: kordoc 4.24.51 -> tools/gen1.py (auto); review=proofread-sonnet; flags=none
\begin{question}{53}
소비자 이론에 관한 설명으로 옳은 것은?
\begin{choices}
  \item 우하향하는 직선의 수요곡선상에서 수요량이 증가할수록 가격탄력성은 감소한다.
  \item 효용함수 $U=X+Y$는 0차 동차 함수이다.
  \item 효용함수 $U=\min\{X , Y\}$에서 $X$재 가격이 상승할 때 $X$재 수요량이 감소하는 것은 대체효과 때문이다.
  \item 소비자가 기펜재와 정상재에 모든 소득을 지출하고 있는 상태에서 기펜재의 가격이 상승하면 정상재에 대한 수요는 증가한다.
  \item 효용함수 $U=X+2Y$이고, 두 재화의 가격이 동일하다면, 소비자는 효용 극대화를 위해 $X$재만을 소비한다.
\end{choices}
\end{question}
\answer{1}
